\section{Decision guide}
\label{sec:decision}

The actionable one page. Starting points below are provisional (from the pre-benchmark, see
\texttt{docs/benchmark/pre\_benchmark.md}); Sec.~\ref{sec:results} confirms or refines them.

\begin{center}
\begin{tabular}{p{4.3cm}p{3.0cm}p{5.6cm}}
\toprule
your data & use & starting parameters \\
\midrule
sparse object, low noise & \textbf{ADMM} & $\kappa \in [0.01, 0.3]$ (raise with noise),
                          $\mu \in [0.5, 1]$ \\
continuous object, noisy & \textbf{ADMM-PnP} & Gaussian/Bilateral, $\sigma \in [25, 50]$,
                          $\rho \ge 0.1$ \\
most versatile / unsure & \textbf{Adam} & TV or SHV, $\lambda \approx 0.1$ (never $>0.3$;
                          $\lambda=0$ only on clean data) \\
uncertainty / MMSE, well-posed & \textbf{MCMC} & TV Bregman, $\sigma \approx 0.02$
                          (deconvolution; on MA-TIRF it does not beat the ridge start) \\
\bottomrule
\end{tabular}
\end{center}

\paragraph{Rules of thumb.}
Noise demands a \emph{stronger} prior --- the best regularisation grows with the noise level.
Match the prior to the structure: sparsity (ADMM) for point-like objects, a denoiser (PnP family)
for continuous ones. Avoid PnP-HQS under noise, and do not spend effort making MCMC win on
MA-TIRF (it structurally cannot). Every parameter in units of $f$ is on a ${\sim}0.02$ scale, not
$1$.
